%% kapa: exam=1 round=23 year=2012 period=1 form=A subject=economics subject_ko=경제원론 seq=18 number=58
%% source: Q-Net 감정평가사 기출문제 (cst00309) artlSeq=5052578 file=raw/1008/5052578/2012년 제23회 감정평가사 제1차시험 기출문제(1교시 A형).pdf
%% conversion: kordoc 4.24.51 -> tools/gen1.py (auto); review=proofread-sonnet; flags=none
\begin{question}{58}
효용극대화를 추구하는 소비자 甲의 효용함수는 $U(x,y)=\min\{x,y\}$이다. 甲의 수요에 관한 설명으로 옳은 것은? (단, 甲은 $X$재와 $Y$재만 소비하고, $x$는 $X$재 소비량, $y$는 $Y$재 소비량을 나타낸다.)
\begin{choices}
  \item 수요의 가격탄력성이 0이다.
  \item 수요의 가격탄력성이 1이다.
  \item 수요의 교차탄력성이 0이다.
  \item 수요의 교차탄력성이 $-1$이다.
  \item 수요의 소득탄력성이 1이다.
\end{choices}
\end{question}
\answer{5}
